\chapter{Color Blindness}
\index{Color Blindness}
\label{sec:Color Blindness}

\section{Overview}
\begin{itemize}[noitemsep]
  \item Color blindness (color vision deficiency) is a reduced ability to distinguish certain colors, most commonly red and green
  \item It affects about 8\% of men and 0.5\% of women worldwide (roughly 1 in 12 men)
  \item Complete color blindness (achromatopsia) is very rare; most people with the condition see colors, just differently
\end{itemize}
\section{Causes}
Most color blindness is inherited in an X-linked pattern, caused by mutations in genes on the X chromosome (OPN1LW and OPN1MW) that encode red and green cone photopigments. Because it is X-linked, it is much more common in males. Acquired color vision deficiency can result from aging, diabetes, glaucoma, retinal disease, certain medications, or chemical exposure.
\section{Risk Factors}
\begin{itemize}[noitemsep]
  \item Male sex (X-linked inheritance)
  \item Family history of color blindness
  \item Eye diseases such as glaucoma, age-related macular degeneration, or diabetic retinopathy
  \item Certain medications and occupational chemical exposures
\end{itemize}
\section{Signs and Symptoms}
\subsection{Common symptoms}
\begin{itemize}[noitemsep]
  \item Difficulty distinguishing red from green, or blue from yellow
  \item Trouble reading color-coded charts, maps, and traffic signals
  \item Colors appear duller or muted
  \item Difficulty identifying ripe fruit or matching clothing
  \item Many people are unaware they have the condition until formally tested
\end{itemize}
\subsection{Emergency warning signs}
\begin{itemize}[noitemsep]
  \item New onset of color vision problems in an adult may signal retinal or optic nerve disease and warrants evaluation
\end{itemize}
\section{Progression and Stages}
Inherited color blindness is present from birth and remains stable throughout life. Acquired deficiencies may progress if the underlying disease progresses. Severity ranges from mild anomalous trichromacy (three cone types present but shifted) to complete dichromacy (one cone type missing).
\section{Complications}
\begin{itemize}[noitemsep]
  \item Difficulty in certain careers (pilot, electrician, police, medicine in some roles)
  \item Safety concerns in tasks relying on color signals
  \item Frustration in daily life and possible anxiety, especially in children
  \item Underlying disease progression in acquired cases
\end{itemize}
\section{Diagnosis}
\begin{itemize}[noitemsep]
  \item Ishihara color plates (most common screening test)
  \item Farnsworth--Munsell 100 hue test to characterize the deficiency
  \item Anomaloscopy to precisely classify the defect
  \item Genetic testing in selected cases
  \item Eye examination to rule out acquired causes when onset is new
\end{itemize}
\section{Prevention}
\begin{itemize}[noitemsep]
  \item Inherited color blindness cannot be prevented
  \item Managing underlying eye diseases may prevent acquired cases
  \item Genetic counseling for families with known inheritance patterns
\end{itemize}
\section{Treatment Options}
\begin{itemize}[noitemsep]
  \item There is no cure for inherited color blindness
  \item Special lenses and filters (including tinted contact lenses) that enhance color discrimination in some people
  \item Gene therapy is under investigation for specific forms
  \item Practical accommodations: labeling, use of apps that identify colors, memorizing signal order
  \item Treatment of the underlying condition in acquired cases
\end{itemize}
\section{Living with It}
\begin{itemize}[noitemsep]
  \item Use smartphone apps and assistive technology for color identification
  \item Arrange clothing and labels using consistent patterns
  \item Inform teachers and employers about appropriate accommodations
  \item Children may benefit from early testing and reassurance
\end{itemize}
\section{Outlook}
Color blindness is a stable, lifelong trait that does not affect general health or life expectancy. With accommodations, most people live full and productive lives with only minor practical limitations.
